\documentclass[12pt]{book}
\usepackage{amsmath}
\usepackage{amssymb}
\usepackage{geometry}
\geometry{margin=1in}

\title{Volume 06: Thermodynamics and Statistical Mechanics \\ Book 02: Transport and Statistical Structure}
\author{James C. Harvey}
\date{December 2025}

\begin{document}
\maketitle
\tableofcontents

\chapter{Transport Coefficients}

\section*{Abstract}
Transport coefficients are derived from matrix contact length scales and flux locking. Viscosity and
diffusivity are geometric consequences rather than empirical primitives.

\section{Viscosity}
\begin{equation}
\mu = \lambda \frac{K_{\text{bulk}} \ell_P^2}{c}.
\end{equation}

\section{Kinematic Viscosity}
Define kinematic viscosity:
\begin{equation}
\nu = \frac{\mu}{\rho}.
\end{equation}
This expresses transport in terms of matrix coupling and density.

\section{Diffusion}
Diffusion arises from flux redistribution across occlusion gradients.

\section{Diffusion Equation}
The diffusion equation is
\begin{equation}
\frac{\partial \rho}{\partial t} = D \nabla^2 \rho,
\end{equation}
where $D$ is the diffusion coefficient derived from matrix contact length scales.

\section{Diffusion Scaling}
An SDT scaling is
\begin{equation}
D \sim c_m \ell_c,
\end{equation}
with $c_m$ the matrix transport speed and $\ell_c$ the contact length.

\chapter{Statistical Structure}

\section*{Abstract}
Statistical structure emerges from the multiplicity of flux configurations. This chapter links
configurational counting to observable transport.

\section{Configuration Counting}
\begin{equation}
\Omega = \exp\left(\frac{S}{k_B}\right),
\end{equation}
which defines the number of accessible flux states.

\section{Fluctuations}
Fluctuations are deviations in local flux occupancy:
\begin{equation}
\langle (\Delta X)^2 \rangle \propto \frac{1}{N_{\text{cells}}},
\end{equation}
where $N_{\text{cells}}$ is the number of participating spations.

\chapter{Thermal Conductivity}

\section*{Abstract}
Thermal conductivity is the macroscopic expression of matrix contact transport.

\section{Conductivity Relation}
\begin{equation}
\kappa_T \sim c_m \ell_c,
\end{equation}
where $c_m$ is matrix transport speed and $\ell_c$ is contact length scale.

\section{Transport Coefficient Link}
The transport coefficients share a common geometric scale:
\begin{equation}
\kappa_T \sim D \sim c_m \ell_c,
\end{equation}
so thermal diffusion and mass diffusion arise from the same matrix-contact mechanism.
\end{document}
